%%%%%%%%%%%%%%%%%%%%%%%%%%%%%%%%%%%%%%%%%%%%%%%%%%%%%%%%%%%%%%%%%%%%%%
% Lớp 12 | Các số đặc trưng ... ghép nhóm | Nguyên hàm
% Dạng 2D4?1-1: Công thức lý thuyết
% Số câu: 29  (sinh tự động bởi tach_cau.py; từ nay thầy sửa trực tiếp file này)
%%%%%%%%%%%%%%%%%%%%%%%%%%%%%%%%%%%%%%%%%%%%%%%%%%%%%%%%%%%%%%%%%%%%%%

%%%%%%%%%%%%%%%% TRẮC NGHIỆM 4 PHƯƠNG ÁN %%%%%%%%%%%%%%%%

% nguồn: Phan1_Lop_12_Dai_so_Chuong_4
\begin{ex}%[2D4N1-1]%[ma:195eb1dd]
	Cho $F(x)$  là một nguyên hàm của $f(x)$ trên $\mathbb{R}$. Chọn khẳng định đúng?
	\choice 
		{\True $\displaystyle\int  f(x) \mathrm{\,d}x=F(x)+C$}
		 {$f'(x)=F(x)$}
		{$\displaystyle\int  F(x) \mathrm{\,d}x=f(x)+C$}
		{$\displaystyle\int F(x) \mathrm{\,d}x=f(x)$} 
	\loigiai{ 
		Ta có $\displaystyle\int f(x)  \mathrm{\,d}x=F(x)+C$.
	} 
\end{ex}

% nguồn: Phan1_Lop_12_Dai_so_Chuong_4
\begin{ex}%[2D4N1-1]%[ma:b104de94]
	Phát biểu nào sau đây là đúng? 
	 \choice
		{$\displaystyle\int \sin x\mathrm{\,d}x = -\sin  x+ C$}
		{$\displaystyle\int \sin x\mathrm{\,d}x =  \sin x+ C$}
		{\True $\displaystyle\int \sin x\mathrm{\,d}x = -\cos  x+ C$}
		{$\displaystyle\int \sin x\mathrm{\,d}x = \cos{x}+  C$}
	\loigiai{ 
		Ta có $\displaystyle\int \sin x\mathrm{\,d}x=-\cos  x+C$.
	} 
\end{ex}

% nguồn: Phan1_Lop_12_Dai_so_Chuong_4
\begin{ex}%[2D4N1-1]%[ma:71bb4202]
	Nếu hàm số $y=f(x)$ có một nguyên hàm trên $\mathbb{R}$ là hàm  số $y=F(x)$ thì
	\choice 
		 {$f'(x)=F(x),\, \forall x\in \mathbb{R}$}
		{$f'(x)=F(x)+C,\, (C\in \mathbb{R}) \, \forall x\in \mathbb{R}$} 
		{\True $F'(x)=f(x), \, \forall  x\in \mathbb{R}$}
		{$F'(x)=f(x)+C,\,(C \in \mathbb{R}\setminus \{0\}),\,\forall x\in  \mathbb{R}$}
	\loigiai{ 
		Ta  có $F'(x)=f(x),\,\forall x\in \mathbb{R}$.
	 }
\end{ex}

% nguồn: Phan1_Lop_12_Dai_so_Chuong_4
\begin{ex}%[2D4N1-1]%[ma:cf4571f6]
	 Cho hàm số $f(x)$ liên tục trên $\mathbb{R}$. Khẳng định nào dưới đây đúng?
	\choice 
		{$\displaystyle\int  2 f(x)\mathrm{\,d}x=\dfrac{1}{2} \displaystyle\int f(x)\mathrm{\,d}x$}
		 {$\displaystyle\int 2 f(x)\mathrm{\,d}x=\displaystyle\int 2 \mathrm{~d} x. \displaystyle\int f(x)\mathrm{\,d}x$}
		{$\displaystyle\int 2 f(x)\mathrm{\,d}x=\displaystyle\int  f(2 x)\mathrm{\,d}x$}
		{\True $\displaystyle\int 2 f(x)\mathrm{\,d}x=2 \displaystyle\int f(x)\mathrm{\,d}x$} 
	 \loigiai{
		Ta có  $\displaystyle\int k f(x)\mathrm{\,d}x=k \displaystyle\int f(x)\mathrm{\,d}x$ với $ k\ne 0 $.
	 }
\end{ex}

% nguồn: Phan1_Lop_12_Dai_so_Chuong_4
\begin{ex}%[2D4N1-1]%[ma:f4f5a1d8]
	Mệnh đề nào sau  đây \textbf{sai}?
	\choice 
		{$\displaystyle\int[f(x)+g(x)]\,\mathrm{d}x=\displaystyle\int f(x)\,\mathrm{d}x+\displaystyle\int g(x)\,\mathrm{d}x$, với  mọi hàm số $f(x)$, $g(x)$ liên tục trên $\mathbb{R}$}
		{$\displaystyle\int f'(x)dx=f(x)+C$ với  mọi hàm số $f(x)$ có đạo hàm trên $\mathbb{R}$}
		{$\displaystyle\int[f(x)-g(x)]\,\mathrm{d}x=\displaystyle\int f(x)\,\mathrm{d}x-\displaystyle\int g(x)\,\mathrm{d}x$, với  mọi hàm số $f(x)$, $g(x)$ liên tục trên $\mathbb{R}$}
		{\True $\displaystyle\int kf(x)\,\mathrm{d}x=k\displaystyle\int f(x)\,\mathrm{d}x$ với mọi hằng số $k$ và với mọi hàm số $f(x)$ liên tục trên $\mathbb{R}$} 
	\loigiai{$\displaystyle\int kf(x)\,\mathrm{d}x=k\displaystyle\int f(x)\,\mathrm{d}x+C$ với mọi hằng số $k\neq 0$  và với mọi hàm số $f(x)$ liên tục trên $\mathbb{R}$.}
\end{ex}

% nguồn: Phan1_Lop_12_Dai_so_Chuong_4
\begin{ex}%[2D4N1-1]%[ma:b4bc1d33]
	Cho hàm số $f(x)$ xác định trên $K$  và $F(x)$ là một nguyên hàm của $f(x)$ trên $K$. Khẳng định nào dưới đây đúng?
	\choice 
		{$f'(x)=F(x)$, $\forall x\in  K$}
		{\True $F'(x)=f(x)$, $\forall  x\in K$}
		 {$f'(x)=F(x)+C$, $\forall x\in K$}
		{$F(x)=f(x)$, $\forall x\in K$} 
	\loigiai{ 
		Ta có $F'(x)=f(x)$, $\forall x\in  K$.
	 }
\end{ex}

% nguồn: Phan1_Lop_12_Dai_so_Chuong_4
\begin{ex}%[2D4N1-1]%[ma:2c67e261]
	Cho hàm số $f(x)$ có đạo hàm là hàm liên tục trên $\mathbb{R}$. Phát biểu  nào sau đây là đúng?
	 \choice
		{$\displaystyle\int  f(x) \mathrm{\,d}x=f'(x)+C$}
		{$\displaystyle\int  f'(x) \mathrm{\,d}x=f(x)$}
		{$\displaystyle\int f(x)  \mathrm{\,d}x=f'(x)$}
		{\True $\displaystyle\int f'(x) \mathrm{\,d}x=f(x)+C$} 
	 \loigiai{
		 Ta có $\displaystyle\int f'(x) \mathrm{\,d}x=f(x)+C$.
	} 
\end{ex}

% nguồn: Phan1_Lop_12_Dai_so_Chuong_4
\begin{ex}%[2D4N1-1]%[ma:6928dca2]
	Hàm số $F(x)$ là một nguyên hàm của hàm số $f(x)$ trên khoảng $K$ nếu 
	 \choice
		{\True  $F'(x)=f(x), \forall x \in K$}
		{$F'(x)=-f(x), \forall  x \in K$}
		{$f'(x)=-F(x), \forall  x \in K$}
		{$f'(x)=F(x), \forall x \in  K$}
	\loigiai{Hàm số $F(x)$ là một nguyên hàm của hàm số  $f(x)$ trên khoảng $K$ nếu $F'(x)=f(x)$, $\forall x \in K$.}
\end{ex}

% nguồn: Phan1_Lop_12_Dai_so_Chuong_4
\begin{ex}%[2D4N1-1]%[ma:19dda676]
	Cho hàm số $f(x)$ liên tục trên đoạn $[0;3]$ và $\displaystyle\int_0^2 f(x)\mathrm{\,d}x =-1$, $\displaystyle\int_0^3 f(x)\mathrm{\,d}x =3$. Tích phân $\displaystyle\int_2^3 f(x)\mathrm{\,d}x$  có giá trị bằng
	 \choice
		{$2$} 
		{$-2$} 
		{$-4$} 
		{\True  $4$}
	\loigiai{ 
		Ta có $\displaystyle\int_0^2 f(x)\mathrm{\,d}x+\displaystyle\int_2^3 f(x)\mathrm{\,d}x=\displaystyle\int_0^3  f(x)\mathrm{\,d}x$.\\
		Từ đó $\displaystyle\int_2^3  f(x)\mathrm{\,d}x=\displaystyle\int_0^3 f(x)\mathrm{\,d}x-\displaystyle\int_0^2 f(x)\mathrm{\,d}x=3-(-1)=4$.
	} 
\end{ex}

% nguồn: Phan1_Lop_12_Dai_so_Chuong_4
\begin{ex}%[2D4N1-1]%[ma:a35ca4e1]
	Trong các khẳng định sau khẳng định nào  \textbf{sai}?
	 \choice
		 {$\displaystyle\int\dfrac{1}{x}\mathrm{\,d}{x}=\ln|x|+C$}
		{$\displaystyle\int{x^\alpha}\mathrm{\,d}{x}=\dfrac{1}{\alpha+1} \cdot x^{\alpha+1}+C;  \quad \alpha \neq -1$}
		 {\True $\displaystyle\int\mathrm{e}^x\mathrm{\,d}{x}=\dfrac{\mathrm{e}^{x+1}}{x+1}+C$}
		 {$\displaystyle\int\sin{x}\mathrm{\,d}{x}=-\cos{x}+C$}
	\loigiai{ 
		$\displaystyle\int  \mathrm{e}^x\mathrm{\,d}{x}=\mathrm{e}^x+C$.
	} 
\end{ex}

% nguồn: Phan1_Lop_12_Dai_so_Chuong_4
\begin{ex}%[2D4N1-1]%[ma:c39e329d]
	Cho số thực $k\ne 0$.  Phát biểu nào sau đây là đúng?
	\choice 
		{$\displaystyle\int k \mathrm{\,d}x =  x+C$}
		 {\True $\displaystyle\int k \mathrm{\,d}x = kx+C$}
		{$\displaystyle\int k \mathrm{\,d}x =  C$}
		{$\displaystyle\int  k \mathrm{\,d}x = kx \cdot C$}
	\loigiai{ 
		Ta  có $\displaystyle\int k\mathrm{\,d}x=kx+C$.
	} 
\end{ex}

% nguồn: Phan1_Lop_12_Dai_so_Chuong_4
\begin{ex}%[2D4N1-1]%[ma:214313c1]
	 Phát biểu nào sau đây là đúng?
	 \choice
		{$\displaystyle\int \cos{x}\mathrm{\,d}x = -\sin{x}+  C$}
		{\True $\displaystyle\int  \cos{x}\mathrm{\,d}x = \sin{x}+ C$}
		{$\displaystyle\int \cos{x}\mathrm{\,d}x =  -\cos{x}+ C$}
		 {$\displaystyle\int \cos{x}\mathrm{\,d}x = \cos{x}+ C$}
	\loigiai{ 
		Ta có  $\displaystyle\int \cos x\mathrm{\,d}x=\sin x+C$.
	 }
\end{ex}

% nguồn: Phan1_Lop_12_Dai_so_Chuong_4
\begin{ex}%[2D4N1-1]%[ma:4fab305d]
	Mệnh đề nào sau  đây {\bf sai}?
	\choice 
		{$\displaystyle\int[f(x)+g(x)]\mathrm{\,d}x  = \int f(x)\mathrm{\,d}x+\int g(x)\mathrm{\,d}x$}
		{$\displaystyle\int[f(x)-g(x)]\mathrm{\,d}x = \int  f(x)\mathrm{\,d}x-\int g(x)\mathrm{\,d}x$}
		{\True $\displaystyle\int k f(x)\mathrm{\,d}x = k\int f(x)\mathrm{\,d}x \,(\forall k \in \mathbb{R})$} 
		{$\displaystyle\int f^m(x)\cdot f'(x)\mathrm{\,d}x  = \dfrac{f^{m+1}(x)}{m+1}+C$; ($m \in \mathbb{R}; m \ne -1$)}
	 \loigiai{
		Mệnh  đề $\displaystyle\int k f(x)\mathrm{\,d}x = k\int f(x)\mathrm{\,d}x$ sai khi $k=0$.
	 }
\end{ex}

% nguồn: Phan1_Lop_12_Dai_so_Chuong_4
\begin{ex}%[2D4N1-1]%[ma:1310e382]
	 Cho hàm số $ y=F(x)$ là một nguyên hàm của hàm số $ y=x^2$. Biểu thức $ F'(25)$ bằng
	 \choice
		{$ 25$} 
		 {\True $ 625$}
		 {$ 5$}
		 {$ 125$}
	 \loigiai{
	 Ta có $ F'(x)=f(x)=x^2$ nên $F'(25)=25^2=625$.}
\end{ex}

% nguồn: Phan1_Lop_12_Dai_so_Chuong_4
\begin{ex}%[2D4N1-1]%[ma:d8f1337e]
	Phát biểu nào sau đây  là đúng?
	 \choice
		{\True $\displaystyle\int F'(x) \mathrm{\,d}  x=F(x)+C$}
		{$\displaystyle\int F(x) \mathrm{\,d}  x=F'(x)+C$}
		{$\displaystyle\int  F(x) \mathrm{\,d} x=F(x)+C$}
		{$\displaystyle\int F'(x)  \mathrm{\,d} x=F'(x)+C$}
	\loigiai{Ta  có $\displaystyle\int F'(x) \mathrm{\,d} x=F(x)+C$}
\end{ex}

% nguồn: Phan1_Lop_12_Dai_so_Chuong_4
\begin{ex}%[2D4N1-1]%[ma:655a7849]
	Hàm số $F(x)=\mathrm{e}^{-2x}$ là nguyên hàm của hàm  số nào sau đây?
	 \choice
		{$f_1(x)=\frac{\mathrm{e}^{-2x}}{-2}$} 
		{$f_2(x)=-\mathrm{e}^{-2x}$} 
		 {$f_3(x)=2\mathrm{e}^{-2x}$}
		{\True  $f_4(x)=-2\mathrm{e}^{-2x}$}
	 \loigiai{Ta có $F'(x)=\big(\mathrm{e}^{-2x}\big)'=-2\mathrm{e}^{-2x}=f_4(x)$.}
\end{ex}

% nguồn: Phan1_Lop_12_Dai_so_Chuong_4
\begin{ex}%[2D4N1-1]%[ma:d812022c]
	Cho hàm  số $f(x)$ có đạo hàm $f'(x)$ liên tục trên $\mathbb{R}$, tìm mệnh đề đúng trong các mệnh đề dưới đây.
	 \choice
		{$\displaystyle\int  f'(x)\mathrm{\,d}x=f(x)$}
		 {$\displaystyle\int f(x)\mathrm{\,d}x=f'(x)$}
		{$\displaystyle\int f(x)\mathrm{\,d}x=f'(x)+C$} 
		{\True $\displaystyle\int f'(x)\mathrm{\,d}x=f(x)+C$} 
	 \loigiai{
	Ta có  $\displaystyle\int f'(x)\mathrm{\,d}x=f(x)+C$.
	} 
\end{ex}

% nguồn: Phan1_Lop_12_Dai_so_Chuong_4
\begin{ex}%[2D4N1-1]%[ma:b2824c80]
	Phát biểu nào sau đây là đúng? 
	\choice 
		{$\displaystyle\int \dfrac{1}{x} \mathrm{\,d}x  =|x|+ C$}
		{\True $\displaystyle\int \dfrac{1}{x} \mathrm{\,d}x =  \ln |x|+C$}
		{$\displaystyle\int \ln{x} \mathrm{\,d}x  =x+ C$}
		{$\displaystyle\int \ln {|x|} \mathrm{\,d}x = \ln{x}+  C$}
	\loigiai{ 
		Ta có $\displaystyle\int \dfrac{1}{x} \mathrm{\,d}x  = \ln |x|+C$.
	} 
\end{ex}

% nguồn: Phan1_Lop_12_Dai_so_Chuong_4
\begin{ex}%[2D4N1-1]%[ma:9d22280e]
	Cho $F\left(x\right)$ là một  nguyên hàm của hàm số $f\left(x\right)$ trên đoạn $\left[a;b\right]$. Mệnh đề nào dưới đây đúng?
	\choice 
		{\True $\displaystyle\int_{a}^{b}f\left(x\right)\,\mathrm{d}x=F\left(b\right)-F\left(a\right)$} 
		 {$\displaystyle\int_{a}^{b}f\left(x\right)\,\mathrm{d}x=F\left(a\right)-F\left(b\right)$}
		{$\displaystyle\int_{a}^{b}f\left(x\right)\,\mathrm{d}x=F\left(b\right)+F\left(a\right)$} 
		{$\displaystyle\int_{a}^{b}  f\left(x\right)\,\mathrm{d}x=-F\left(b\right)-F\left(a\right)$}
	\loigiai{Ta  có $\displaystyle\int_{a}^{b}f\left(x\right)\,\mathrm{d}x=F\left(b\right)-F\left(a\right)$.}
\end{ex}

% nguồn: Phan1_Lop_12_Dai_so_Chuong_4
\begin{ex}%[2D4N1-1]%[ma:7ac7a791]
	Phát biểu nào sau  đây là đúng?
	\choice 
		{$\displaystyle\int 0 \mathrm{\,d}x =  -x+C$}
		{$\displaystyle\int 0 \mathrm{\,d}x  =x+ C$}
		 {$\displaystyle\int 0 \mathrm{\,d}x = Cx,\, C\in \mathbb{R}\setminus \left\lbrace 0\right\rbrace $}
		 {\True $\displaystyle\int 0 \mathrm{\,d}x = C$}
	\loigiai{ 
		Ta  có $\displaystyle\int 0\mathrm{\,d}x=0x+C=C$.
	} 
\end{ex}

% nguồn: Phan1_Lop_12_Dai_so_Chuong_4
\begin{ex}%[2D4N1-1]%[ma:02038089]
	Phát biểu  nào sau đây \textbf{sai}?
	 \choice
		{$\displaystyle\int\big(f(x)+g(x)\big) \mathrm{\,d}x=\displaystyle\int f(x) \mathrm{\,d} x+\displaystyle\int  g(x) \mathrm{\,d} x$}
		{$2 \displaystyle\int f(x)  \mathrm{\,d}x=\displaystyle\int 2 f(x) \mathrm{\,d} x$}
		{\True $\displaystyle\int x f(x) \mathrm{\,d}x=x \displaystyle\int f(x) \mathrm{\,d}  x$}
		{$\displaystyle\int\big(f(x)-g(x)\big)  \mathrm{\,d}x=\displaystyle\int f(x) \mathrm{\,d} x-\displaystyle\int g(x) \mathrm{\,d} x$}
	 \loigiai{
		Phát biểu \lq\lq$\displaystyle\int x f(x) \mathrm{\,d}x=x \displaystyle\int f(x)  \mathrm{\,d}x$\rq\rq\ là phát biểu sai.
	 }
\end{ex}

% nguồn: Phan1_Lop_12_Dai_so_Chuong_4
\begin{ex}%[2D4N1-1]%[ma:cd508571]
	Cho hàm số $F(x)$ là một nguyên hàm của hàm số  $f(x)$ trên $\mathbb{R}$. Mệnh đề sau, mệnh đề nào sai.
	 \choice
		{$\displaystyle\int f(x)dx = F(x)+C$} 
		{$\left(\displaystyle\int f(x)dx\right)'  = f(x)$}
		{\True $\left(\displaystyle\int f(x)dx\right)'  = f(x)+C$}
		{$\left(\displaystyle\int f(x)dx\right)' =  F'(x)$}
	 \loigiai{
		$\left(\displaystyle\int f(x)dx\right)' = F(x)+C$ 
	} 
\end{ex}

% nguồn: Phan1_Lop_12_Dai_so_Chuong_4
\begin{ex}%[2D4N1-1]%[ma:c16829fd]
	Cho $\displaystyle\int5^{x} \mathrm{~d} x=F(x)+C$. Khẳng định nào dưới đây \textbf{đúng}? 
	\choice 
		{$F'(x)=5^x\cdot\ln  5$}
		 {$F'(x)=-5^x$}
		{\True $F'(x)=5^x$} 
		{$F'(x)=5^x+C$} 
	 \loigiai{
	Lấy đạo hàm hai  vế ta được $F'(x)=5^x$.}
\end{ex}

% nguồn: Phan1_Lop_12_Dai_so_Chuong_4
\begin{ex}%[2D4N1-1]%[ma:1141c261]
	Hàm số  $F(x)$ là một nguyên hàm của hàm số $f(x)$ trên khoảng $K$ nếu
	\choice 
		{\True $F'(x)=f(x)$, $\forall x \in  K$}
		{$F'(x)=-f(x)$,  $\forall x \in K$}
		{$f'(x)=-F(x)$, $\forall  x \in K$}
		 {$f'(x)=F(x)$, $\forall x \in K$}
	\loigiai{Hàm số $F(x)$ là một nguyên hàm của hàm số $f(x)$ trên khoảng  $K$ nếu $F'(x)=f(x)$, $\forall x \in K$.}
\end{ex}

% nguồn: Phan1_Lop_12_Dai_so_Chuong_4
\begin{ex}%[2D4N1-1]%[ma:090d6bfd]
	Trong  các khẳng định sau, khẳng định nào \textbf{sai}?
	 \choice
		 {$\displaystyle\int 0 \mathrm{\,d}x=C$}
		 {$\displaystyle\int \dfrac{1}{x} \mathrm{\,d}x=\ln |x|+C$}
		{\True  $\displaystyle\int x^a\mathrm{\,d}x=\dfrac{x^{a+1}}{a+1}+C$}
		{$\displaystyle\int  \mathrm{\,d}x=x+C$}
	\loigiai{$\displaystyle\int x^a \mathrm{\,d}x=\dfrac{x^{a+1}}{a+1}+C$ không  đúng với trường hợp $a=-1$.}
\end{ex}

% nguồn: Phan1_Lop_12_Dai_so_Chuong_4
\begin{ex}%[2D4H1-1]%[ma:8b468863]
	Hàm số $F(x)=x^2+\sin x$ là một nguyên  hàm của hàm số
	\choice 
		 {$f(x)=2x-\cos x$}
		{$f(x)=\dfrac{1}{3} x^3-\cos  x$}
		 {$f(x)=\dfrac{1}{3} x^3+\cos x$}
		 {\True $f(x)=2x+\cos x$}
	 \loigiai{
		Ta có $f(x)=F'(x)=2x+\cos x$. 
	} 
\end{ex}

% nguồn: Phan1_Lop_12_Dai_so_Chuong_4
\begin{ex}%[2D4H1-1]%[ma:19b82035]
	Biết $F(x)=x^2$ là một nguyên hàm của hàm số $f(x)$ trên $\mathbb{R}$. Tích  phân $\displaystyle\int_{1}^2{\left[2+f(x)\right] \mathrm{\,d}x}$ có giá trị bằng
	 \choice
		 {$\dfrac{13}{3}$}
		 {$3$}
		 {\True $5$}
		{$\dfrac{7}{3}$} 
	 \loigiai{
		Ta có $\displaystyle\int_{1}^2{\left[2+f(x)\right]\mathrm{\,d}x=\displaystyle\int_{1}^2{2}\mathrm{\,d}x}+\displaystyle\int_{1}^2{f(x)}\mathrm{\,d}x$.\\ 
		Mà $\displaystyle\int_{1}^2{f(x)}\mathrm{\,d}x=F(2)-F(1)=2^2-1^2=4-1=3$.\\ 
	Vậy $\displaystyle\int_{1}^2{\left[2+f(x)\right]  \mathrm{\,d}x}=2 \cdot (2-1)+3=2+3=5$.}
\end{ex}

% nguồn: Phan1_Lop_12_Dai_so_Chuong_4
\begin{ex}%[2D4H1-1]%[ma:e525ce21]
	Cho $\displaystyle\int_{1}^{2}\left[4f(x)-2x\right]\mathrm{\,d}x=1$. Khi đó $\displaystyle\int_{1}^{2}f(x)\mathrm{\,d}x$ bằng 
	\choice 
		 {\True $1$}
		 {$-3$}
		{$3$} 
		 {$-1$}
	 \loigiai{
		Ta  có $\displaystyle\int_{1}^{2}\left[4f(x)-2x\right]\mathrm{\,d}x=1$. Suy ra $4\displaystyle\int_{1}^{2}f(x)\mathrm{\,d}x-\displaystyle\int_{1}^{2}2x\mathrm{\,d}x=1\Leftrightarrow\displaystyle\int_{1}^{2}f(x)\mathrm{\,d}x=1$.
	 }
\end{ex}

% nguồn: Phan1_Lop_12_Dai_so_Chuong_4
\begin{ex}%[2D4H1-1]%[ma:92d77e54]
	Cho hàm số $F(x)$ là một nguyên hàm của hàm số $f(x)$ trên $K$.  Trong các mệnh đề sau, mệnh đề nào là \textbf{sai}?
	 \choice
		 {$\displaystyle\int f(x)\mathrm{\,d}x=F(x)+C$}
		{$\left( \displaystyle\int  f(x)\mathrm{\,d}x\right)'=f(x)$}
		{\True $\left(  \displaystyle\int f(x)\mathrm{\,d}x\right)'=f'(x)$}
		{$\left( \displaystyle\int  f(x)\mathrm{\,d}x\right)'=F'(x)$}
	\loigiai{Vì $\left( \displaystyle\int f(x)\mathrm{\,d}x\right)'=f(x)$ nên mệnh  đề $\left( \displaystyle\int f(x)\mathrm{\,d}x\right)'=f'(x)$ sai.}
\end{ex}
